\paragraph{Stages and tiles.}\label{app:two}
A \emph{stage of scale $t$}, for fixed $0<t\leq 1$, has a fixed collection
of groups.  Each group has an input list and an output list of the same
length $2^{th+O(1)}$, and the lists have $O(n)$ positions in total.  A
stage specifies $n$ matched input--output pairs, with both positions in
the same group and with each position used at most once on its side.
We call each matched pair a \emph{tile}.  A permutation on $n$ positions
is a stage of scale one with a single group.


The construction groups positions by successive colors. A recursive
permutation network uses a related decomposition \cite{Benes}. Here
the control must also compute the route in linear word time on every
input matching.

\par\Needspace{4\baselineskip}
\begin{lemma}[Subdivision of a stage]\label{lem:subdivision}
Fix $0<t\leq 1$ and an integer $d\geq2$, and put $K=18(d-1)+1$.
For all sufficiently large $h$, a stage of scale $t$ can be replaced by
$K$ successive stages of scale $t/d$.  Each parent tile is routed from
its input to its output along one tile of every child stage, within the
original parent group.  At each of the $K+1$ boundaries between moves,
including the two ends, different parent tiles occupy distinct positions.
The child groups and the identifications between successive position
lists are fixed independently of the matching.  Given the matching, the
control computes the child tiles in $O(n)$ operations on $O(h)$-bit
words, with $O(n)$ words of advice and working storage.  The constants
may depend on $t,d$ and on the size bounds for the parent groups.
\end{lemma}

\begin{proof}
Put $D=2^{\lfloor th/(9d)\rfloor}$ and $L_b=D^9$.  Inside one parent
group, begin with the given pair of lists.  To split a pair of lists of
length $m_1$, partition each list into blocks of $L_b$ positions, padding
the last block.  Color every matched pair with one of $D$ colors so that
at either endpoint at most $4D^8$ pairs of its color use its block.
For a greedy assignment, a block with $4D^8$ previous pairs of one color
forbids that color.  Each endpoint block forbids at most
$L_b/(4D^8)=D/4$ colors, so a color is always available at both ends.
We give the time bound for the assignment below.

In each block reserve $4D^8$ positions for every color, in a fixed order.
One entering stage sends each occupied input position to the next free
position of its color in the same block.  One exiting stage sends the
corresponding color positions on the output side to their final
positions.  Allocate the output-side positions in the same processing
order, using a separate counter for each block and color.  Both stages
use groups with $4L_b$ positions on each side after padding.  Between
them there remains, for each color, a matching between its positions
across all blocks.  Its two lists each have length
\[
        m_2=4D^8\lceil m_1/L_b\rceil.
\]
The fixed relabeling by color and block identifies each position of an
entering or exiting stage with a unique position of the intervening lists.

Repeat this operation on each remaining list pair for $q=9(d-1)$ rounds.
Starting from $m_0=2^{th+O(1)}=\Theta(D^{9d})$, the length after $j$
rounds is $\Theta(D^{9d-j})$ for $0\leq j\leq q$.  Indeed, before
the last split the length is at least a constant times $D^{10}$, and
each split of a list of length $m_1\geq L_b$ gives
$4m_1/D\leq m_2\leq8m_1/D$.  Thus every final list has length
$\Theta(D^9)=2^{th/d+O(1)}$ and supports one central stage.  Reading
the $q$ entering stages, the central stage, and the $q$ exiting stages
in reverse nesting order gives the asserted $K$ moves.  In one round
the $D$ separated lists together have length at most $8m_1$.  Since the
number of rounds is fixed, each phase inside the parent group has
$O(m_0)$ positions and can be padded to a common length
$\Theta(m_0)$.  Each matched pair uses distinct endpoints at each step
by its allocated positions and the inductive matching property.

It remains to assign colors in linear time.  Let $\ell=\lceil Ch\rceil$
for an absolute constant $C$.  Hardwire $D^4$ strings of $\ell$ colors
so that, for every $S\subseteq\{1,\ldots,D\}$ of size at most $D/4$,
fewer than $D$ strings have at least half their entries in $S$.
To obtain the strings, sample their entries independently and uniformly.
For fixed $S$ the probability for one string is $\exp(-\Omega(\ell))$.
The probability of failure for some $S$ and some $D$ strings is at most
$2^D\binom{D^4}{D}\exp(-\Omega(\ell D))<1$ for sufficiently large $C$,
since $\log D\leq h$.

For each block store a bit mask for each string.  A bit is set when the
corresponding color is forbidden in this block.  When a color first
becomes forbidden, set its bits using the list of its occurrences across
the strings.  Initialization and all such updates cost $O(D^4\ell)$
word operations per block.  For a matched pair with within-block
positions $u,v$, try the string indexed by
$(u\bmod D^2,v\bmod D^2)$.  Take the union of its two endpoint masks.
If there is a free bit, use its color.  This takes constant time on
$O(h)$-bit words.  For example, divide a mask into a constant number of
chunks of at most $h/2$ bits and use a shared table of $O(2^{h/2})$
words to find the first free bit in a chunk.

If both masks together cover the string, scan all $D$ colors to find a
common available one.  Charge this scan to an endpoint block whose mask
covers at least half the string.  Its final set of forbidden colors has
size at most $D/4$.  By the choice of strings, fewer than $D$ strings
can receive charges at that block.  A fixed string is tried there at
most $L_b/D^2=D^7$ times, since it fixes the residue of the position
on that side and the matching uses each position at most once.  Thus all
scans charged to a block cost at most $D\cdot D\cdot D^7=L_b$ word
operations.  Also $D^4\ell=O(L_b)$ for fixed $t,d$ and sufficiently
large $h$.  Summing over the blocks in the fixed number of rounds proves
the claimed bounds.  Fixing the order of the pairs and taking the first
available color in either search makes the routing deterministic.

The scans, shifts, comparisons and table reads used here fit
Definition~\ref{def:protocol}. If a word operation needs multiplication
by an advice-fixed constant, split the operands into a bounded number of
pieces of at most $h/3$ bits. A shared product table of $O(2^{2h/3})$
words and a bounded number of additions and shifts suffice.
\end{proof}
